../cictikz/src/cictikz/data/tex/cictikz_preamble.tex
% The ring oscillator SUN_PLL_ROSC, drawn from the xschem netlist in
% sun_pll_sky130nm. A NAND and eight inverters make nine inversions, so
% the loop oscillates whenever PWRUP_1V8 is high. The ring gates run on
% VDD_ROSC - that supply IS the control node, which is why it is the one
% thing drawn in red. The level shifter takes two adjacent taps (N2, N1)
% back up to the AVDD domain, and an output inverter buffers CK.
\begin{circuitikz}[american, thick, transform shape, circuit ee IEC]
../cictikz/src/cictikz/data/tex/cictikz_lib.tex

% Enable NAND and the eight ring inverters (hand-rolled gates)
\draw (0,0) \cicNand{nd};
\draw (nd_out)
  -- ++(0.3,0) \cicInv{I0} -- ++(0.3,0) \cicInv{I1}
  -- ++(0.3,0) \cicInv{I2} -- ++(0.3,0) \cicInv{I3}
  -- ++(0.3,0) \cicInv{I4} -- ++(0.3,0) \cicInv{I5}
  -- ++(0.3,0) \cicInv{I6} -- ++(0.3,0) \cicInv{I7} coordinate (ringEnd);

% Close the ring: last output back to the lower NAND input
\draw (ringEnd) -- ++(0.4,0) -- ++(0,-1.5) -| (nd_in2);

% Enable port
\draw (nd_in1) -- ++(-0.6,0) \portIn{PWRUP\_1V8};

% The control domain: everything inside the red box runs on VDD_ROSC
\draw[red, dashed, rounded corners] (-0.9,-1.9) rectangle (13.1,0.9);
\draw[red] (0.2,1.6) to [short,o-] ++(0.6,0) -- ++(0,-0.7)
  node [anchor=south east, red] {VDD\_ROSC};

% Level shifter on the AVDD domain, fed from taps N2 and N1
\draw (11.6,2.2) rectangle ++(1.8,1.6);
\node[anchor=center] at (12.5,3.0) {LS};
\coordinate (lsInTop) at (11.6,3.4);
\coordinate (lsInBot) at (11.6,2.6);
\coordinate (lsOut) at (13.4,3.0);

% Taps: N2 is the output of I5, N1 the output of I6
\draw (I5_out) ++(0.15,0) to [short,*-] ++(0,3.4) node [anchor=south east] {N2} -| (lsInTop);
\draw (I6_out) ++(0.15,0) to [short,*-] ++(0,2.6) node [anchor=south east] {N1} -| (lsInBot);

% Output buffer
\draw (lsOut) \cicInv{I8};
\draw (I8_out) \portOut{CK};

\end{circuitikz}
\end{document}
